\section{Two further benchmarks sharpen the evaluation contract}
\label{exec14:l7-frontier}

Two primary studies extend the comparison in different directions. Wang's
\emph{CausalBench} groups four perspectives on each scenario across text,
mathematics and code. Babakov, Reiter and Bugarín's \emph{CausalGraphBench}
standardizes descriptions for reconstructing published graphs. The names
refer to distinct artifacts; Wang's language benchmark is also distinct from
the biological perturbation benchmark discussed elsewhere. Their value to
the proposal comes from their task definitions and exposed limitations,
not from treating every reported score as a common causal capability.
\citep{exec14_l7_causalbench,exec14_l7_causalgraphbench}

\subsection{Four questions make a stricter outcome, not a proof of understanding}

CausalBench asks cause-to-effect and effect-to-cause questions, each with and
without intervention. Its construction starts from manually expanded
scenarios, uses GPT-4 Turbo for scale, and describes a causal inference engine
followed by human review when answers disagree. The four-question group
score requires all four answers to be correct. Table 2 reports 61.4\% group
correctness and 88.7\% overall correctness for GPT-4 Turbo on the mathematical
tasks. These are historical author-reported outcomes. They do not measure
the effect of language assistance on a banking decision.
\citep{exec14_l7_causalbench}

The metric has a useful, exact interpretation. If \(S_j\) is correctness
on perspective \(j\), \(p_j=\Pr(S_j=1)\), and
\(q=\Pr(S_1=S_2=S_3=S_4=1)\), then
\[
       \max\{0,\textstyle\sum_{j=1}^4p_j-3\}
       \leq q\leq\min_j p_j.
\]
The lower bound follows by applying the union bound to the four error
events; the upper bound follows because joint success implies each
individual success. Neither argument assumes independent answers. Using
the reported GPT-4 Turbo mathematical margins, 78.6\%, 88.3\%, 91.7\%
and 96.0\%, gives 54.6\%--78.6\%. The reported 61.4\% lies inside that
range. The group score measures their joint behavior, which the mean
alone conceals. It does not establish how the correct answers were produced.
Even four independent fair binary guesses are jointly correct with probability
\(1/16\); dependence or answer imbalance changes that probability.
\citep{exec14_l7_frontier_exact}

The quality-control description also illustrates the adjudication problem.
Agreement between a generated answer and an engine is accepted; disagreement
triggers manual analysis. Shared graph-extraction, target or identification
errors can survive agreement. In particular, the paper's generic suggestion
of an instrumental-variable or front-door route for hidden confounding
requires the corresponding assumptions and information; hidden confounding
alone does not supply either route. The project should retain the four-view
rubric while independently checking formal targets and a probability sample
of agreed as well as disputed labels.

Direct dataset adoption remains conditional. The inspected proceedings text
describes approximately 10,000 text, 2,000 mathematics and 2,000 code scenarios
with four questions each, while its introduction says more than 60,000
problems. Exact counts and split/version identity require the original
dataset bytes. Its CLADDER bibliography attribution also differs from the
Jin et al. original inspected here. These discrepancies should be reconciled
before claiming inherited provenance. The described engine is not an
implementation verified in this workspace. The proposed disposition is to
adapt and test the rubric, while deferring imported labels and code pending
primary artifact, license and oracle inspection.
\citep{exec14_l7_causalbench}

\subsection{Graph reconstruction and a recall probe have limited targets}

CausalGraphBench selects 35 Bayesian-network graphs: 14 from repositories
and 21 from papers. The repository graphs have median 42 nodes, compared
with 14 for the paper graphs. Runnable files and the feasibility of extracting
diagrams affect selection. GPT-4o helps extract graph ideas, domains and
clarified variable descriptions from source material; the evaluation gives
methods standardized descriptions. The resulting task is reconstruction
from semantic information. It is not discovery from an observed joint law,
and the authors explicitly leave data-driven and hybrid comparisons to
future work. Each imported Bayesian network would still need review of
its causal interpretation for the proposed banking use.
\citep{exec14_l7_causalgraphbench}

Method comparisons must retain their size scopes. Pair and triplet querying
are limited to graphs with at most ten nodes, and the iterative efficient
method to fifty. On the small-graph GPT-4o comparison, Table 2 reports
edge-normalized structural Hamming distances of 0.65 for the single-query
baseline, 1.67 for pairs and 2.02 for triplets. These adverse results make
the simple baseline essential. They do not support a compute-matched
universal ranking. The study also reports manual repair of malformed
outputs and generation-token limitations; a banking workflow must charge
that intervention to its actual operational outcome.

The metric convention is material: reversed directions count as two errors,
and normalization uses the reference edge count. A graph with a different
normalization or reversal convention cannot be compared numerically without
recalculation. Nor does this directed-reference score automatically fit an
equivalence-class output. Graph validity, unresolved orientation and
downstream effect support retain their own evaluation contracts.
\citep{exec14_l7_causalgraphbench}

The contamination probe asks a model to recall nodes from a named source,
then recover edges for high-recall cases. Its stated rule uses less than
15\% deviation in node count and recall above 0.85. This is a useful
diagnostic for recognizable public graphs, but no calibrated sensitivity
or specificity for training exposure is supplied. Failure to recall a source
does not prove absence from training. There is also an internal inconsistency
to resolve before copying the flags: the Llama-3.3 \emph{asiam} row has nine
generated nodes against seven true nodes and is treated as known, although
the count deviation is \(2/7\), or about 28.6\%, above the stated limit.
The printed record supports reporting this discrepancy; it does not justify
inventing a different threshold or correcting unpublished raw results.
\citep{exec14_l7_causalgraphbench,exec14_l7_frontier_exact}

The proposal therefore adopts the idea of standardized descriptions and
source-recall diagnostics as development tools. Importing graph files,
prompts or scores awaits artifact and rights inspection, resolved count
rules, checked variable meanings and an explicit causal reference.
Node paraphrasing can be an additional diagnostic only when the meaning
is preserved. Neither the count probe nor paraphrasing supplies the
exposure bound in the preceding mixture analysis. This restricted reuse
route retains valuable evidence while keeping the next prototype testable.

\label{exec14:l7-frontier-end}
